\section{AC feasibility and the meaning of angle limits}
\label{sec:ac}

The resistive result transfers to an AC model with zero reactive injections.
The transfer requires a consistent real angle at every bus: small principal
angle differences on individual lines need not fit together around a cycle.
We first encode this consistency without transcendental constraints. The
encoding works for every line limit strictly below $\pi$, including limits
larger than $\pi/2$.

\subsection{The AC input model}

Write $\mathrm{i}^2=-1$. On a simple undirected graph, line $\{i,j\}$ has
series admittance $y_{ij}=g_{ij}+\mathrm{i}b_{ij}$, where
$g_{ij}\in\Q_{\ge0}$, $b_{ij}\in\Q$, and $y_{ij}\ne0$. Bus $i$ may have a
shunt admittance $s_i=h_i+\mathrm{i}t_i$, with $h_i\in\Q_{\ge0}$ and
$t_i\in\Q$; absence of a shunt means $h_i=t_i=0$. These are reciprocal
series lines without transformers or phase shifters. A line-charging shunt
at a line endpoint may be included in that bus's shunt sum.
Let $U_i=v_i\exp(\mathrm{i}\theta_i)$, with $v_i>0$ and
$\theta_i\in\R$. With injection positive into the network, the equations are
\begin{equation}
 I_i=s_iU_i+\sum_{j:\{i,j\}\in E}y_{ij}(U_i-U_j),\qquad
 P_i+\mathrm{i}Q_i=U_i\overline{I_i}.
 \label{eq:complex-power}
\end{equation}
The input gives positive rational magnitude bounds $\ell_i\le v_i\le u_i$
and finite rational intervals $[p_i^-,p_i^+]$, $[q_i^-,q_i^+]$ for $P_i,Q_i$.
Singleton intervals are allowed independently for these three quantities.
Every line also has a rational number $c_{ij}\in(-1,1]$, representing
$\gamma_{ij}=\arccos c_{ij}\in[0,\pi)$, and the constraint
\begin{equation}
 |\theta_i-\theta_j|\le\gamma_{ij}.
 \label{eq:real-angle-limit}
\end{equation}
The problem $\ACPF$ asks for $v\in\R_{>0}^N$ and $\theta\in\R^N$ satisfying
these equations and bounds. The rational input contains $c_{ij}$, not an
approximation to $\gamma_{ij}$. Angles may be shifted by a common constant
on each connected component; an optional reference angle zero per component
only fixes this freedom.

For later use, write $U_i=e_i+\mathrm{i}f_i$ and introduce a positive magnitude
variable $r_i$. Define
\begin{equation}
 H_{ij}=e_ie_j+f_if_j,\qquad D_{ij}=f_ie_j-e_if_j.
 \label{eq:dot-cross}
\end{equation}
Thus $D_{ij}=\operatorname{Im}(U_i\overline{U_j})=-D_{ji}$.
Expanding~\eqref{eq:complex-power} gives the polynomial expressions
\begin{align}
 P_i&=h_ir_i^2+\sum_j\bigl[g_{ij}(r_i^2-H_{ij})-b_{ij}D_{ij}\bigr],
 \label{eq:rect-P}\\
 Q_i&=-t_ir_i^2+\sum_j\bigl[-b_{ij}(r_i^2-H_{ij})-g_{ij}D_{ij}\bigr],
 \label{eq:rect-Q}\\
 r_i^2&=e_i^2+f_i^2,\qquad \ell_i\le r_i\le u_i.
 \label{eq:rect-magnitude}
\end{align}
Here and below, line sums run over neighbors of $i$. All these constraints
have rational coefficients and degree at most two.

\subsection{A polynomial encoding of real angle lifts}

For two nonzero phasors with non-antipodal directions, let
$\delta_{ij}\in(-\pi,\pi)$ be the principal argument of
$U_i\overline{U_j}$. The principal-angle limit
$|\delta_{ij}|\le\gamma_{ij}$ is exactly
\begin{equation}
 H_{ij}\ge c_{ij}r_ir_j.
 \label{eq:cosine-test}
\end{equation}
Indeed, $H_{ij}/(r_ir_j)=\cos\delta_{ij}$, and cosine decreases on $[0,\pi]$.
Since $c_{ij}>-1$, antipodal directions are excluded. The positive variables
$r_i$ make squaring unnecessary, so the same test covers negative cosines.
Condition~\eqref{eq:cosine-test} alone does not encode~\eqref{eq:real-angle-limit}.

\begin{lemma}[Short-arc crossing count]
\label{lem:crossing}
Let $U_i,U_j\ne0$ have non-antipodal directions, and traverse the short
oriented arc from $U_j/|U_j|$ to $U_i/|U_i|$. Its signed crossing count of
the positive real ray, counting arrival from the open lower half-plane as
$+1$ and departure into that half-plane as $-1$, is
\begin{equation}
 k_{ij}=\begin{cases}
 1,& f_j<0\le f_i\ \text{and }D_{ij}>0,\\
 -1,& f_i<0\le f_j\ \text{and }D_{ij}<0,\\
 0,&\text{otherwise}.
 \end{cases}
 \label{eq:crossing-rule}
\end{equation}
In particular $k_{ji}=-k_{ij}$. If $\alpha_i,\alpha_j\in[0,2\pi)$ are the
arguments of the endpoints, then
\begin{equation}
 \delta_{ij}=\alpha_i-\alpha_j+2\pi k_{ij}.
 \label{eq:crossing-identity}
\end{equation}
\end{lemma}
\begin{proof}
Put $a=\alpha_i-\alpha_j\in(-2\pi,2\pi)$.
If $a< -\pi$, then $\delta_{ij}=a+2\pi\in(0,\pi)$.
Necessarily $\alpha_j>\pi$ and $\alpha_i<\pi$, so
$f_j<0\le f_i$, and $D_{ij}=r_ir_j\sin\delta_{ij}>0$.
Conversely, $f_j<0\le f_i$ implies $\alpha_j\in(\pi,2\pi)$ and
$\alpha_i\in[0,\pi]$, hence $a<0$. Within $(-2\pi,0)$ the additional
condition $\sin a>0$ holds exactly when $a< -\pi$.
Thus the first case of~\eqref{eq:crossing-rule} is precisely $a< -\pi$.
Reversing $i,j$ proves that its second case is precisely $a>\pi$.
The remaining case has $|a|<\pi$, including $a=0$, and needs no wrap.
The equalities $a=\pm\pi$ are excluded by the non-antipodal hypothesis.
This proves~\eqref{eq:crossing-identity} and reversal symmetry. It also
checks the axis convention: argument zero belongs to the closed upper
half-plane; argument $\pi$ does too, but the determinant sign prevents
counting a crossing of the negative real ray.
\end{proof}

Fix one orientation $j\to i$ for every line. For any directed traversal of
a cycle $C$, write $\varepsilon_{C,ij}=1$ if it traverses that line in its
chosen orientation and $-1$ if it traverses it in reverse. Telescoping the
$\alpha$ terms in~\eqref{eq:crossing-identity} gives
\begin{equation}
 \sum_{\{i,j\}\in C}\varepsilon_{C,ij}\delta_{ij}
 =2\pi\sum_{\{i,j\}\in C}\varepsilon_{C,ij}k_{ij}.
 \label{eq:cycle-winding}
\end{equation}
The integer on the right after division by $2\pi$ is the winding number.
Cycle-angle recovery and winding representations are established in
\citet[Theorem 2]{FarivarLow2013} and \citet{JafarpourEtAl2022}.
Here the selected short differences must sum to zero over the reals,
rather than merely modulo $2\pi$.

\begin{lemma}[Real-angle lift criterion]
\label{lem:lift}
Suppose~\eqref{eq:cosine-test} holds on every line. There are real angles
$\theta_i$ representing the phasors and satisfying~\eqref{eq:real-angle-limit}
if and only if
\begin{equation}
 \sum_{\{i,j\}\in C}\varepsilon_{C,ij}k_{ij}=0
 \label{eq:zero-winding}
\end{equation}
for each fundamental cycle of a spanning forest of $G$.
Equivalently, there exist real vertex variables $a_i$ satisfying
\begin{equation}
 a_{\mathrm{root}}=0\quad\text{in each component},\qquad
 a_i-a_j=k_{ij}\quad\text{on each oriented line }j\to i.
 \label{eq:winding-potentials}
\end{equation}
Whenever they exist, these root-normalized variables are unique integers.
\end{lemma}
\begin{proof}
For a real lift within the limits, both $\theta_i-\theta_j$ and
$\delta_{ij}$ belong to $(-\pi,\pi)$ and represent the same argument.
They are therefore equal, and their sum around every cycle is zero.
Equation~\eqref{eq:cycle-winding} proves necessity.
Conversely choose one argument at a root of each tree of the forest and
extend $\theta$ along its edges by $\theta_i-\theta_j=\delta_{ij}$.
These angles represent the given phasors. Each non-tree edge and its
unique forest path form a fundamental cycle. Its zero sum in
\eqref{eq:cycle-winding} forces the same difference equation on that edge.
Every line difference is thus its principal difference and satisfies its limit.

For the equivalent formulation, propagate $a_i-a_j=k_{ij}$ along the
forest from $a_{\mathrm{root}}=0$. This uniquely fixes every $a_i$ as a
sum of signed integer crossing counts, including $a_i=0$ at isolated
vertices. Each remaining edge equation holds exactly when the crossing
sum on its fundamental cycle is zero. Conversely, vertex differences
telescope to zero on every cycle. In particular, no integer quantifier
is needed: integrality follows from the edge equations and the root
normalization. The angles $\theta_i=\alpha_i+2\pi a_i$ then give an
explicit real lift by~\eqref{eq:crossing-identity}. For an empty graph,
all conditions are vacuous.
\end{proof}

\begin{theorem}
\label{thm:ac-membership}
The problem $\ACPF$ belongs to $\ER$ for rational cosine limits
$c_{ij}\in(-1,1]$. Its existential encoding uses $O(|N|+|E|)$ real
variables, atomic predicates, and monomial occurrences, all of degree
at most two, and has polynomial total bit length.
\end{theorem}
\begin{proof}
Use variables $e_i,f_i,r_i$ subject to~\eqref{eq:rect-magnitude}, the injection
intervals applied to~\eqref{eq:rect-P}--\eqref{eq:rect-Q}, and
\eqref{eq:cosine-test}. For each oriented line introduce a real $k_{ij}$ and
encode~\eqref{eq:crossing-rule} by a Boolean combination of polynomial
comparisons: either its first condition and $k_{ij}=1$, or its second
condition and $k_{ij}=-1$, or the negations of both conditions and $k_{ij}=0$.
Add the vertex variables and equations~\eqref{eq:winding-potentials}.
Lemmas~\ref{lem:crossing}--\ref{lem:lift} prove equivalence to real-angle
feasibility. There are four real variables per vertex, one per line,
one normalization per component, and one difference equation per line.
The crossing rule has constant size per line, and the expanded injection
expressions have $O(|N|+|E|)$ monomial occurrences in total. Thus the
stated linear counts hold, including disconnected and empty graphs.
Variable indices and rational input coefficients still contribute to the
total bit length; we claim polynomial, not linear, total bit length.
Clearing denominators preserves degree at most two and has polynomial
bit cost. All quantified variables are real; the crossing rule and
root-normalized difference equations enforce their stated discrete values.
\end{proof}

The vertex representation is the elementary potential form of a zero
cycle-sum condition. It avoids listing long fundamental cycles in the
formula; the crossing rule supplies its connection to rational phasor
coordinates.

\subsection{Zero reactive injection and the hardness transfer}

For purely resistive lines with $g_{ij}>0$ and no shunts, the polar form of
\eqref{eq:rect-P}--\eqref{eq:rect-Q} is
\begin{align}
 P_i&=\sum_j g_{ij}\bigl(v_i^2-v_iv_j\cos(\theta_i-\theta_j)\bigr),
 \label{eq:resistive-ac-P}\\
 Q_i&=-\sum_j g_{ij}v_iv_j\sin(\theta_i-\theta_j).
 \label{eq:resistive-ac-Q}
\end{align}

\begin{lemma}[Equal angles under a real lift]
\label{lem:equal-angles}
In a purely resistive, shunt-free network, suppose $v_i>0$, $Q_i=0$ at
every bus, and $|\theta_i-\theta_j|<\pi$ on every line. Then angles are
constant on each connected component. The active injections are exactly
$\power_i(v)$ from~\eqref{eq:rpf}. Conversely, every positive resistive
voltage profile extends to an AC profile with these injections and $Q=0$
by taking all angles zero.
\end{lemma}
\begin{proof}
Multiply~\eqref{eq:resistive-ac-Q} by $-\theta_i$ and sum over buses.
Pairing the two contributions of each line yields
\begin{equation}
 0=-\sum_i\theta_i Q_i
 =\sum_{\{i,j\}\in E}g_{ij}v_iv_j
 (\theta_i-\theta_j)\sin(\theta_i-\theta_j).
 \label{eq:angle-energy}
\end{equation}
Each summand is nonnegative, and $t\sin t>0$ for $0<|t|<\pi$.
Since $g_{ij}v_iv_j>0$, every line difference is zero. Substitution in
\eqref{eq:resistive-ac-P} gives~\eqref{eq:rpf}. The converse follows directly
from~\eqref{eq:resistive-ac-P}--\eqref{eq:resistive-ac-Q}.
\end{proof}

The sine balance in~\eqref{eq:resistive-ac-Q} is also the equilibrium equation
of a weighted oscillator network; see \citet{DorflerEtAl2013} for this
connection and cohesive-angle analysis. The elementary identity
\eqref{eq:angle-energy} supplies the exact statement needed here, with real
bus angles and line differences strictly below $\pi$. It does not require
strictly positive cosine or a stability assumption.

\begin{theorem}
\label{thm:ac-complete}
The problem $\ACPF$ is $\ER$-complete. Hardness holds on simple graphs of
maximum degree three, with $b_{ij}=h_i=t_i=0$, $q_i^-=q_i^+=0$ at every bus,
$c_{ij}=0$ on every line, and the fixed finite conductance, magnitude, and
active-injection data of Theorem~\ref{thm:rpf-complete}.
\end{theorem}
\begin{proof}
Take the resistive instance of Theorem~\ref{thm:rpf-complete}, retain its
graph and data, and add the stated AC data. The cosine input zero represents
the real-angle limit $\pi/2$. Lemma~\ref{lem:equal-angles} proves equivalence
in both directions, and Theorem~\ref{thm:ac-membership} proves membership.
\end{proof}

Thus $\ACPF\in\NP$ would imply $\NP=\ER$. This conclusion concerns the
specified AC model with variable magnitudes and resistive lines. It does not
classify the lossless, fixed-magnitude model, whose nonlinear variables are
angles alone.

\begin{corollary}[One-sided reactive intervals]
\label{cor:one-sided-reactive}
Theorem~\ref{thm:ac-complete} remains true if every reactive-injection
interval $[0,0]$ is replaced by the positive-width interval $[0,1]$.
More generally, in a purely resistive, shunt-free network, requiring all
reactive injections to be nonnegative, or all to be nonpositive, forces
each of them to be zero.
\end{corollary}
\begin{proof}
Antisymmetry of the sine terms in~\eqref{eq:resistive-ac-Q} gives
$\sum_i Q_i=0$. Numbers of one common weak sign with zero sum all vanish.
Conversely, zero belongs to $[0,1]$, so the feasible set is unchanged.
\end{proof}

The same replacement applies to the two restricted AC formulations below.
This removes the need to specify singleton reactive intervals, while the
network identity still enforces exact zero reactive injections. It is a
one-sided saturation argument, not a robustness assertion for symmetric
intervals $[-\varepsilon,\varepsilon]$.

\subsection{Principal angles: a counterexample and a positive-window result}
\label{subsec:principal}

Define \emph{principal-angle AC feasibility} by
\eqref{eq:rect-P}--\eqref{eq:cosine-test} and the injection intervals, without
the cycle equations~\eqref{eq:zero-winding}. This is a polynomial feasibility
problem and hence belongs to $\ER$. The distinction from $\ACPF$ is material.

\begin{proposition}[Winding prevents an unrestricted transfer]
\label{prop:winding-counterexample}
On a unit-conductance $n$-cycle, $n\ge3$, the phasors
$U_k=\exp(2\pi\mathrm{i}k/n)$ have unit magnitudes and satisfy
\begin{equation}
 Q_k=0,\qquad P_k=2\bigl(1-\cos(2\pi/n)\bigr)>0.
 \label{eq:cycle-profile}
\end{equation}
Their principal line differences have absolute value $2\pi/n$ and winding
number one. Consequently, for every fixed rational $c\in(-1,1)$ there is a
rational-data instance with line cosine limit $c$ that is principal-angle
AC feasible but whose corresponding resistive instance is infeasible.
\end{proposition}
\begin{proof}
At each bus, the two sine terms in~\eqref{eq:resistive-ac-Q} cancel and the
cosine terms coincide, giving~\eqref{eq:cycle-profile}. The oriented sum
around the cycle is $2\pi$, so Lemma~\ref{lem:lift} excludes a real lift
with these principal differences. Given $c<1$, choose $n\ge3$ large enough
that $2\pi/n\le\arccos c$. Put unit magnitude bounds and zero reactive
injections at all buses, and choose rational numbers $0<a<P_k<b$ as active
injection bounds $[a,b]$. Such numbers exist by density of $\Q$; the powers
in~\eqref{eq:cycle-profile} themselves need not be rational. The phasors
give principal-angle feasibility, whereas every resistive voltage is one
and its injection is zero, outside $[a,b]$.
For an explicit instance with singleton rational injections, take $n=4$,
$c=0$, and $[p_k^-,p_k^+]=[2,2]$.
\end{proof}

This example disproves an equal-angle implication from fixed positive
principal windows; it does not assert a complexity classification for every
such fixed-window class. At the other endpoint, the bound strictly below
$\pi$ in Lemma~\ref{lem:equal-angles} is sharp: a single unit-conductance
line with unit magnitudes and real angles $0,\pi$ has $Q=0$ and active
injection $2$ at both buses.

A positive principal window can nevertheless enforce the needed real lift
when its width is allowed to depend on the graph size.

\begin{lemma}[Small cycle budgets exclude winding]
\label{lem:small-cycle-budget}
Suppose principal angle limits $\gamma_{ij}<\pi$ satisfy
$\sum_{\{i,j\}\in C}\gamma_{ij}<2\pi$ for every fundamental cycle $C$
of some spanning forest. Then every phasor profile satisfying those limits
has a real lift satisfying the same limits.
\end{lemma}
\begin{proof}
The principal difference sum on $C$ is a multiple of $2\pi$ by
\eqref{eq:cycle-winding}. Its absolute value is at most the stated sum of
limits and therefore strictly less than $2\pi$. It must vanish.
Lemma~\ref{lem:lift} applies. In particular, every forest has this property
without any additional restriction on its limits below $\pi$.
\end{proof}

\begin{corollary}[Strictly positive principal windows]
\label{cor:principal-hardness}
Principal-angle AC feasibility is $\ER$-complete even on simple graphs of
maximum degree three with purely resistive lines, no shunts, zero reactive
injections, the fixed numerical data of Theorem~\ref{thm:rpf-complete}, and
the common cosine limit
\begin{equation}
 c_n=1-\frac{1}{n^2},\qquad n=|N|\ge2.
 \label{eq:shrinking-window}
\end{equation}
The angle window is strictly positive. Its cosine has $O(\log n)$ bits;
unlike the other data, these cosines do not belong to a fixed finite set.
\end{corollary}
\begin{proof}
For $t\in[0,\pi]$, concavity of sine on $[0,\pi/2]$ gives
$\sin(t/2)\ge t/\pi$, and therefore
$1-\cos t=2\sin^2(t/2)\ge2t^2/\pi^2$.
Writing $\gamma_n=\arccos c_n$, we obtain
$0<\gamma_n\le\pi/(\sqrt{2}n)$.
Every simple cycle has at most $n$ edges, so the sum of its limits is at
most $\pi/\sqrt{2}<2\pi$. Lemma~\ref{lem:small-cycle-budget} supplies a
real lift, and Lemma~\ref{lem:equal-angles} proves equivalence to the
resistive instance. Take that instance from Theorem~\ref{thm:rpf-complete}.
If needed, add isolated buses with magnitude bounds $[1/2,2]$ and active
injection bounds $[-85,85]$ to ensure $n\ge2$. These buses are feasible
and use the original finite alphabet. Computing~\eqref{eq:shrinking-window} takes
polynomial time. Membership follows from the rectangular formulation.
\end{proof}

\subsection{A rectangular model with bus-angle boxes}
\label{subsec:angle-box}

There is also a hardness transfer using only linear angle restrictions in
rectangular coordinates. Choose a reference bus in each connected component
and require
\begin{equation}
 e_i\ge f_i,\qquad e_i\ge-f_i\quad(i\in N),\qquad
 f_{\mathrm{ref}}=0,\quad e_{\mathrm{ref}}>0.
 \label{eq:angle-box}
\end{equation}
The positive lower magnitude bounds imply $e_i>0$, so these conditions are
exactly the angle box $\theta_i\in[-\pi/4,\pi/4]$ relative to a reference
angle zero. Their real line differences lie in $[-\pi/2,\pi/2]$.

\begin{corollary}
\label{cor:box-hardness}
The rectangular AC feasibility problem with~\eqref{eq:angle-box} is
$\ER$-complete under the same graph, admittance, magnitude, and injection
restrictions as Theorem~\ref{thm:ac-complete}, without separate line-angle
constraints. In this restricted class every feasible profile has $f_i=0$
and $e_i=v_i$ at every bus, and its feasible set is exactly the resistive
voltage set in these coordinates.
\end{corollary}
\begin{proof}
The formulation is polynomial. The box supplies real angles with the
required line differences, so Lemma~\ref{lem:equal-angles} makes them
constant on each component. Its reference angle makes that constant zero.
Conversely every feasible resistive profile with $e_i=v_i$, $f_i=0$
satisfies the box and all AC constraints.
\end{proof}

The three transfers have different input conventions. Real line limits
permit the fixed choice $c_{ij}=0$ and require the winding equations for
membership. Bus-angle boxes give a direct polynomial encoding with fixed
coefficients. Positive principal-only windows give a direct polynomial
encoding with the size-dependent cosine~\eqref{eq:shrinking-window}.
None of these equivalences assumes that small numerical reactive residuals
are exactly zero.
